% -------------------------------------------------------------------------
% ÜBERHOLT (31.08.2026). Maßgeblich ist die eingereichte main.tex. Zahlen,
% Zeilennummern und Angaben zum Overleaf-Stand in dieser Datei stammen aus
% der Zeit VOR der Crosswalk-Korrektur vom 29.08.2026 (35 statt 36 Stadt-
% teile, vollständiger Neulauf) und wurden bewusst NICHT nachgezogen: die
% Datei ist als datierter Entwurf Teil der Arbeitsdokumentation.
% -------------------------------------------------------------------------
% =========================================================================
% KAPITEL 5 -- GEKUERZTE FASSUNG, ALLE VIER UNTERABSCHNITTE
% 25.08.2026. Ersetzt den Fliesstext von 5.1 bis 5.4 vollstaendig.
% Gleitobjekte bleiben, wo sie stehen; sieben kommen hinzu
% (docs/kapitel5_gekuerzt_listings.tex).
%
% UMFANG
%   vorher  3.806 Woerter Fliesstext, 10 Quelltexte
%   nachher rund 2.900 Woerter, 15 Quelltexte
%   5.1  600 -> 465   5.2  468 -> 415
%   5.3  979 -> 730   5.4 1.759 -> 1.290 (17 -> 13 Absaetze)
%
% WAS BEIM KUERZEN GEFALLEN IST
%   - Einzelwerte, wo eine Spanne genuegt (ACS-Trefferquoten, Brandfaelle
%     je Fold, Trainings- und Testzeilen je Fold)
%   - das jeweils dritte Argument, wo zwei tragen
%   - Bezeichner, die der Quelltext ohnehin zeigt
%   - Ergebnisse, die nach Kapitel 7 gehoeren (Fold-4-R2, Faktor 32,
%     Trefferquote 0,806 gegen 0,584, Poisson gegen NegBin als RMSE-Paar
%     an ZWEI Stellen)
%
% WAS BEWUSST STEHEN GEBLIEBEN IST
%   Auflage A (Absatz 3 von 5.4), Auflage B (woertlich, Absatz 1 von 5.4),
%   Auflage C (Absatz 7), Auflage D (woertlich, Absatz 11) samt Vorbehalt
%   R-2, R-5 (Absatz 6), R-14 (Absatz 11), Serialisierung (Absatz 8).
%
% ALLE FRUEHER GEMELDETEN FEHLER SIND EINGEARBEITET: das fehlende
% Leerzeichen, die halbierte darf/kann-Zweiteilung, der Bezug auf das
% Expose statt auf Unterfrage 2, E(i,t) statt E(i), die beiden woertlichen
% Dopplungen mit 4.2, die Standardisierung, die Herleitung von "keine
% fehlenden Werte".
% =========================================================================


% ####################### 5.1 ############################################

\subsection{Datenauswahl und Bereinigung}
\label{sec:auswahl-bereinigung}

Der Rohbestand aus Abschnitt~\ref{sec:datengrundlage} umfasst 720.258
Einsatzmeldungen aus 41 Analysis Neighborhoods. Nicht alles davon darf und
nicht alles kann in die Analyse eingehen. Die erste Frage betrifft die
Zulässigkeit einer Größe und folgt dem Grundsatz aus
Abschnitt~\ref{sec:modellbewertung}, dass ein Prädiktor zum Prognosezeitpunkt
verfügbar gewesen sein muss; die zweite grenzt den Gegenstand zeitlich und
räumlich ab und stützt sich auf die Qualitätsbefunde aus
Abschnitt~\ref{sec:eda}. Beide Entscheidungen sind nicht über den
Verarbeitungscode verteilt, sondern als benannte Konstanten hinterlegt
(Quelltext~\ref{lst:konstanten}); jede ist damit an einem Ort nachlesbar und an
einem Ort änderbar. Am Ende steht der Rahmen, in dem gerechnet wird: 35
Stadtteile über 132 Monate.

In die Analysedatensätze gelangen ausschließlich Größen, die zum Zeitpunkt des
Alarms bereits feststanden --- das no-time-machine
requirement.\footcite[559]{Kaufman2011} Damit sind zehn Felder des
Feldverzeichnisses gesperrt, darunter Sachschaden, eingesetzte Kräfte und
Antwortzeit (Quelltext~\ref{lst:gesperrt}). Sie sind nicht nur unzulässig,
sondern gefährlich: Sie hängen eng mit der Zielgröße zusammen, und ein Modell
mit ihnen erreichte hohe Gütewerte, ohne eine Prognose zu leisten. Der
Ausschluss geschieht deshalb beim Laden und nicht erst beim Modellieren ---
was nicht in die Datei gelangt, kann später nicht versehentlich in ein Modell
geraten.\footcite[560]{Kaufman2011} Dass keines dieser Felder den Weg in die
erzeugten Dateien findet, wird nicht behauptet, sondern geprüft
(Quelltext~\ref{lst:test-leakage}).

Bereinigt wird nur, was nachweislich Artefakt ist. 269 doppelte Einsatznummern
werden entfernt (Abschnitt~\ref{sec:eda}); Ausrückzeiten außerhalb von 0 bis 60
Minuten und Baujahre außerhalb von 1800 bis 2025 werden als fehlend markiert
statt gelöscht. Beide Grenzen sind gesetzte Plausibilitätsannahmen und folgen
nicht aus den Daten. An der Zielgröße selbst wird nichts gekappt oder getrimmt;
Bereinigung betrifft ausschließlich Platzhalterwerte der Quellen.

Sechs Stadtteile scheiden ganz aus; sie sind in Abschnitt~\ref{sec:eda}
benannt --- drei Parkgebiete ohne nennenswerte Wohnbevölkerung und drei ohne
durchgängige Abdeckung durch die ACS-Jahrgänge. Entscheidend ist die Form
dieses Ausschlusses: Entfernt werden vollständige Analyseeinheiten, nicht
einzelne Zeilen. Ein zeilenweises Verwerfen erzeugte ein unbalanciertes Panel,
in dem ein Stadtteil mitten in der Zeitreihe hinzuträte und Summen im
Testfenster allein durch diesen Zutritt sprängen. Ein rechteckiges Panel ist
Voraussetzung dafür, dass die Folds aus Abschnitt~\ref{sec:rahmen-baselines}
miteinander vergleichbar sind.

Der Analysezeitraum ist auf Januar 2015 bis Dezember 2025 festgelegt, also auf
132 Monate. Der Beginn ist keine Setzung, sondern eine Folge der
Verfügbarkeitsregel: Für Einsätze vor 2015 wäre der jüngste zulässige
ACS-Jahrgang der von 2009, und der führt die Tabelle nicht, aus der die
Akademikerquote stammt (Abschnitt~\ref{sec:eda}). Das Ende ist das letzte
vollständige Kalenderjahr. Beide Grenzen sind fest hinterlegt und werden nicht
aus den Daten abgeleitet, weil die Analyse sonst mit jedem neuen Abzug wanderte
und ein angebrochener Randmonat unbemerkt ins Testfenster geriete.


% ####################### 5.2 ############################################

\subsection{Zusammenführung der Datenquellen}
\label{sec:integration}

Die sieben Quellen aus Abschnitt~\ref{sec:datengrundlage} liegen in drei
Raumbezügen und zwei Zeitrastern vor. Sie werden auf der Ebene des einzelnen
Einsatzes zusammengeführt und nicht auf der späteren Analyseeinheit: Jeder
Einsatz trägt seinen eigenen Zeitpunkt, und nur auf dieser Ebene lässt sich ihm
derjenige Stand der zeitabhängigen Merkmale zuordnen, der damals verfügbar war.
Verdichtet wird erst danach (Abschnitt~\ref{sec:merkmale}); die Reihenfolge
folgt damit der Verarbeitung und nicht der Aufgabenliste des Vorgehensmodells.
Das Ergebnis ist eine Tabelle mit 719.989 Zeilen und 50 Spalten, eine je
Einsatz.

Die sozioökonomischen Merkmale werden in zwei Stufen angefügt. Räumlich ordnet
ein Crosswalk jeden Census Tract einem Stadtteil zu; dabei werden Zählgrößen
summiert, Medianwerte dagegen bevölkerungsgewichtet gemittelt, weil sich ein
Median nicht addieren lässt. Zeitlich gilt die Bedingung
$\text{ACS-Jahrgang} \le \text{Einsatzjahr} - 1$, sodass einem Einsatz nur ein
Jahrgang zugeordnet wird, der zu seinem Zeitpunkt bereits veröffentlicht war
(Quelltext~\ref{lst:acs-versatz}). Vollständig ist die räumliche Zuordnung
nicht: Der Crosswalk beruht auf den Tract-Grenzen des Zensus 2020, sodass die
Trefferquote je nach Jahrgang zwischen 63,1 und 99,2\,\% liegt.

Die Kriminalitätsbelastung stammt aus zwei Quellen, deren Schnitt im Januar
2018 liegt. Die Altquelle führt keine Stadtteilspalte, dafür Koordinaten; ihre
Meldungen werden über einen räumlichen Verbund derselben Geometrie zugeordnet,
die auch für die baulichen Merkmale gilt --- sonst bezögen sich Kriminalitäts-
und Baumerkmale desselben Stadtteils auf unterschiedliche Flächen
(Quelltext~\ref{lst:spatialjoin}). Geschnitten wird am Beginn der
Nachfolgequelle, damit kein Monat aus zwei Erfassungssystemen gespeist wird.

Die baulichen Merkmale liegen je Parzelle vor und werden über den
Flächenmittelpunkt zugeordnet; das gelingt für 154.544 Parzellen und damit für
99,5\,\% des Verzeichnisses. Der Bestand ist ein Abzug des Jahres 2020 und der
einzige verfügbare Jahrgang. Die baulichen Merkmale sind deshalb je Stadtteil
über den gesamten Analysezeitraum konstant (Abschnitt~\ref{sec:eda}).

Alle drei Zusammenführungen sind nach demselben Muster abgesichert: Vor jedem
Verbund wird die Eindeutigkeit des Schlüssels geprüft, nach jedem die Zeilenzahl
gegen den Erwartungswert gehalten, und bei Abweichung bricht die Aufbereitung
ab (Quelltext~\ref{lst:joinhygiene}) --- ebenso, wenn die Zuordnungsquote des
räumlichen Verbunds unter 90\,\% fällt. Das ist keine Sorgfaltsroutine: Ist ein
Schlüssel mehrdeutig, entsteht ein kartesisches Produkt, und jeder betroffene
Einsatz geht mehrfach in die spätere Aggregation ein.


% ####################### 5.3 ############################################

\subsection{Konstruktion der Merkmale und Zielgrößen}
\label{sec:merkmale}

Die Einsatztabelle aus Abschnitt~\ref{sec:integration} führt 719.989 Zeilen,
eine je Ereignis. Auf dieser Ebene stehen Merkmale und Zielgröße nicht auf
derselben Beobachtung: Die Strukturmerkmale beschreiben einen Stadtteil, die
Zielgröße dessen Einsatzlast. Die Verdichtung auf Stadtteil und Monat stellt
beides in dieselbe Zeile und bildet dabei die Größen, mit denen gerechnet wird.
Eine Verdichtung auf Census Tract und Jahr scheidet aus, weil die Einsatzdaten
keine Tract-Zuordnung tragen und die Saisonalität dabei verloren ginge.
Tabelle~\ref{tab:codebook} führt die entstehenden Größen mit Skalenniveau,
Einheit und Wertebereich; die folgenden Absätze begründen ihre Bildung.

Verdichtet wird auf ein vollständiges Raster aus 35 Stadtteilen und 132
Monaten: Auch ein Monat ohne Einsatz erhält eine Zeile, und zwar mit dem Wert
null statt gar nicht. Fehlende Merkmalswerte werden ausschließlich vorwärts
fortgeschrieben; ein Rückwärtsfüllen schlösse eine Lücke mit einem später
erhobenen Wert und spielte damit genau die Information ein, die zum
Prognosezeitpunkt nicht vorlag (Quelltext~\ref{lst:raster}). Am Ende stehen
4.620 Zeilen für die Menge der Einsatzlast --- 35 mal 132, lückenlos --- und
4.619 für ihre Zusammensetzung; dem Strukturdatensatz fehlt der eine Monat ohne
Einsatz.

Sechs der zehn Strukturmerkmale sind Anteilswerte aus je zwei Zählgrößen
derselben Erhebung. Ein Nenner von null ergibt dabei einen fehlenden Wert und
nicht null, damit eine fehlende Bezugsgröße sichtbar bleibt
(Quelltext~\ref{lst:quoten}). Ein Nenner ist begründungspflichtig: Der
Altbauanteil bezieht sich auf die Parzellen mit bekanntem Baujahr, weil er
sonst dort am stärksten zu niedrig ausfiele, wo das Verzeichnis am
lückenhaftesten ist (Abschnitt~\ref{sec:eda}).

Das kriminalitätsbezogene Merkmal ist kein Rohwert, sondern ein
Standortquotient: die Deliktrate eines Stadtteils, gemessen an der Deliktrate
der Gesamtstadt im selben Monat.

\begin{equation*}
    \mathrm{Rate}(i,t) = \frac{D(i,t)}{E(i,t)}
    \qquad
    \mathrm{Index}(i,t) = \frac{\mathrm{Rate}(i,t)}{\mathrm{Rate}(\mathrm{Stadt},t)}
\end{equation*}

Dabei ist $D(i,t)$ die Zahl der Delikte im Stadtteil $i$ über die zwölf Monate
bis einschließlich $t-1$ und $E(i,t)$ seine Einwohnerzahl im zugeordneten
Zensusjahrgang; ein Wert von eins bedeutet eine Belastung wie im
Stadtdurchschnitt desselben Monats. Die Form des Standortquotienten stammt aus
der Regionalwissenschaft und wird von Brantingham und Brantingham auf
Kriminalität übertragen.\footcite[268-269]{Brantingham1997} Deren Fassung
verwendet allerdings die Gesamtzahl der Delikte als Bezugsgröße und misst
damit, worauf ein Gebiet relativ spezialisiert ist; hier ist die Bezugsgröße
die Wohnbevölkerung, gemessen wird die relative Belastung je Einwohner.

Zwei Eigenschaften dieser Konstruktion sind Entscheidungen. Die relative Form
ist die Antwort auf den Quellenwechsel aus Abschnitt~\ref{sec:integration}: Ein
Wechsel des Erfassungssystems verschiebt das stadtweite Niveau annähernd
multiplikativ und kürzt sich deshalb weitgehend heraus. Was bleibt, ist eine
Verschiebung in der Zusammensetzung der erfassten Delikte, die einzelne
Stadtteile stärker trifft als andere. Die zweite Eigenschaft ist das Fenster:
Es endet strikt im Vormonat, weil eine im selben Monat gemessene Kriminalität
die Einsätze dieses Monats erklären statt vorhersagen würde
(Quelltext~\ref{lst:krimindex}). Der Index geht logarithmiert ein, sodass der
Stadtdurchschnitt bei null liegt.

Die Wohnbevölkerung geht als eigenes Merkmal ein, ebenfalls logarithmiert. Sie
wirkt auf die beiden Mengenzielgrößen gegenläufig: Ein größerer Stadtteil hat
mehr Einsätze, aber nicht zwangsläufig mehr Einsätze je Einwohner. Ohne diese
Größenkontrolle bildete ein Modell auf der absoluten Zielgröße vor allem die
Einwohnerzahl ab und nicht die Struktur, nach der gefragt ist. Im
Poisson-Modell des Abschnitts~\ref{sec:rahmen-baselines} tritt sie zusätzlich
als Offset auf.

Der Kalendermonat wird nicht als Zahl von eins bis zwölf geführt, sondern als
Sinus- und Kosinusterm. Als Zahl hätten Dezember und Januar den Abstand elf,
obwohl sie aufeinanderfolgen. Die beiden Terme sind die einzigen Merkmale, die
ausschließlich innerhalb der Stadtteile variieren (Abschnitt~\ref{sec:eda}).

Drei Verlaufsmerkmale --- Vormonat, Vorjahresmonat und Mittel der drei
Vormonate --- werden strikt rückwärtsgerichtet und je Stadtteil gebildet; damit
sie schon für den ersten Analysemonat definiert sind, beginnt die Aggregation
zwölf Monate früher, ohne dass diese Vorlaufmonate eigene Zeilen bilden. In die
Modelle gehen sie dennoch nicht ein: Unter der Aufteilung des
Abschnitts~\ref{sec:rahmen-baselines} wäre der Vormonatswert die eigene
Vergangenheit eines zurückgehaltenen Stadtteils. Sie verbleiben im Datensatz,
weil die Beschreibung in Abschnitt~\ref{sec:datengrundlage} sie benötigt.

Drei Zielgrößen stehen nebeneinander (Abschnitt~\ref{sec:zielgroessen}): die
Zahl der Einsätze je Stadtteilmonat, dieselbe Zahl auf je 1.000 Einwohner
normiert und die überwiegende Einsatzart. Die vier Anteile der NFIRS-Gruppen
sind dabei Zwischengröße und keine eigenen Zielgrößen; aus ihnen entsteht die
überwiegende Einsatzart als der größte der vier. Lage und Streuung der beiden
Mengenzielgrößen stehen in Abschnitt~\ref{sec:eda}.

Dass die Einsatzart auf Stadtteilebene und nicht je Einsatz bestimmt wird, ist
die folgenreichste Festlegung dieses Abschnitts. Innerhalb eines
Stadtteilmonats tragen alle Einsätze dieselben Strukturmerkmale; die 350.481
Einsätze des Analysezeitraums verteilen sich deshalb auf nur 4.619 verschiedene
Merkmalsprofile. Ein Modell, das jedem Profil die dort häufigste Einsatzart
zuwiese, träfe in 49,9\,\% der Fälle zu, die häufigste Klasse ohne jedes
Merkmal zu nennen in 48,2\,\% --- der gesamte Spielraum beträgt 1,7
Prozentpunkte. Auf Stadtteilebene ist die Frage dagegen beantwortbar. Die
Zielgröße bleibt dabei eine echte Klasse und kein an einem Schwellwert
geteilter stetiger Anteil; damit entfällt die Begründungslast einer solchen
Teilung, deren Wirkung Altman und Royston mit dem Verlust eines Drittels der
Beobachtungen beziffern.\footcite{Altman2006}


% ####################### 5.4 ############################################

\subsection{Analyserahmen und Baselines}
\label{sec:rahmen-baselines}

Die Einsatzhäufigkeit unterscheidet sich weit stärker zwischen den Stadtteilen
als innerhalb eines Stadtteils über die Zeit; auf die räumliche Dimension
entfallen 92,5\,\% der Varianz (Abschnitt~\ref{sec:eda}). Ein Modell, das einen
Stadtteil bereits aus dem Training kennt, muss dessen Niveau deshalb nicht mehr
erklären. Die Aufteilung in Trainings- und Testmengen wird daher nicht entlang
der Zeit, sondern entlang der Stadtteile gezogen; Ziel der Validierung ist die
Generalisierung auf unbekannte Stadtteile. Sie entsteht in der Aufbereitung und
wird als Spalten \texttt{fold} und \texttt{ist\_holdout} in beide
Analysedateien geschrieben, ist also eine Eigenschaft des Datensatzes und nicht
der Algorithmen --- das Vorgehensmodell führt sie folgerichtig unter der
Auswahl der Daten.\footcite[23-26]{Chapman2000}

Geprüft wird, indem ganze Stadtteile zurückgehalten werden. Die 29
Entwicklungsstadtteile verteilen sich auf fünf Folds mit sechs, sechs, sechs,
sechs und fünf Stadtteilen; jeder ist genau einmal Testfall, und zwar mit allen
132 Monaten (Abbildung~\ref{abb:foldstruktur}, Quelltext~\ref{lst:test-folds}).
Ein Zeitschnitt prüfte die Forschungsfrage nicht: Dort steht jeder Stadtteil in
Training und Test, und das Modell kennt sein Niveau bereits --- der
Stadtteilmittelwert allein erreicht ein Bestimmtheitsmaß von 0,888
(Abschnitt~\ref{sec:eda}). Für genau diesen Fall empfehlen Roberts et al.\ eine
räumlich geblockte Aufteilung, wenn auf neue Orte vorhergesagt werden
soll.\footcite[925]{Roberts2017} Was der Rahmen dadurch nicht mehr misst, ist
die Fortschreibung eines bereits bekannten Stadtteils in die Zukunft.

Damit weicht die Untersuchung von der angemeldeten Fassung ab. Das Exposé sah
für die zweite Unterfrage eine zeitreihengerechte Kreuzvalidierung und ein
Testfenster der letzten zwölf Monate vor. Der Wechsel ist am 28.~Juli 2026
entschieden worden, nachdem die Varianzzerlegung des Bestands vorlag; er ist
die Folge eines Befundes über die Daten und keine Anpassung an ein Ergebnis.
Die Abwägung beider Rahmen steht in Kapitel~\ref{sec:diskussion}.

Welcher Stadtteil in welchen Fold fällt, ist nicht zufällig. Sortiert wird
zuerst nach der Zahl brand-dominierter Monate, bei Gleichstand nach der
Wohnbevölkerung; die so geordneten Stadtteile werden reihum auf sechs Gruppen
verteilt, von denen die erste das Hold-out bildet (Quelltext~\ref{lst:folds}).
Das erste Kriterium ist nötig, weil die Hälfte der 70 brand-dominierten Monate
auf einen einzigen Stadtteil entfällt: Ohne Stratifizierung blieb ein Fold
regelmäßig ganz ohne Brand-Testfall, und das Macro-F1 mittelte über eine
Klasse, die dort nicht vorkam. Das zweite verhindert, dass die Streuung
zwischen den Folds ein Größeneffekt ist. Ein Leckageproblem entsteht nicht:
Festgelegt wird allein, welche Stadtteile gemeinsam getestet werden.

Sechs Stadtteile bleiben während der gesamten Entwicklung unberührt und werden
erst für die Schlussbewertung geöffnet. Sie sind die erste Gruppe derselben
Ordnung und deshalb kein verkleinertes Abbild der Stadt: Auf sie entfallen 37
der 70 brand-dominierten Monate, ihr Brandanteil liegt bei 4,67\,\% gegenüber
0,86\,\% in den Entwicklungsstadtteilen. Ergebnisse auf dem Hold-out sind daher
nicht als Durchschnittsleistung über die Stadt zu lesen, sondern als Probe
unter erschwerten Bedingungen.

Eine einzelne Aufteilung von 29 Stadtteilen auf fünf Folds ist eine
willkürliche Gruppierung. Alle Referenzwerte entstehen deshalb aus 50 Läufen,
fünf Folds in zehn Wiederholungen; berichtet wird die Streuung dieser zehn
Wiederholungsmittel und nicht die über alle 50 Läufe, der dieselben 29
Stadtteile in zehn Anordnungen zugrunde lägen
(Quelltext~\ref{lst:streuung}). Auch die Zahl der Zeilen überzeichnet, wie viel
unabhängige Information vorliegt. Weil das Raster rechteckig ist, gilt die
Näherung für gleich große Gruppen hier genau: Aus einer Intraklassenkorrelation
von 0,926 und 132 Monaten je Stadtteil folgt ein Designeffekt von 122 und damit
eine effektive Stichprobe von rund 38 statt 4.620
Beobachtungen.\footcite[206]{Killip2004} Die Aufteilung nach Stadtteilen trägt
dem Rechnung, indem sie den Stadtteil und nicht die Zeile als Einheit
behandelt.

Dass alle Verfahren dieselben Merkmale und dieselbe Aufteilung sehen, ist keine
Zusicherung, sondern eine Eigenschaft des Aufbaus. Die Fold-Zuteilung entsteht
in einem einzigen Aufruf und wird von dort in beide Analysedateien geschrieben;
Mengen- und Strukturstrang können konstruktiv keine unterschiedlichen
Aufteilungen sehen. Die Merkmalsliste steht genau einmal, in
\texttt{prep/config.py}, und wird von jedem Modellskript von dort gelesen
(Quelltext~\ref{lst:merkmalsliste}). Eine automatisierte Prüfung baut die
Aufteilung aus den Dateien nach und vergleicht sie mit den gespeicherten
Spalten; ohne sie liefen bei einer veralteten Aufteilung alle Verfahren auf
falschen, aber untereinander identischen Mengen: Der Vergleich bliebe fair und
das Ergebnis wäre trotzdem falsch.

Zuletzt werden beide Dateien auf modelltaugliche Typen gebracht: Merkmale als
\texttt{float64}, Schlüssel und Zählgrößen als \texttt{int64}, keine fehlenden
Werte, weil das Raster vollständig ist (Quelltext~\ref{lst:dtypes}). Das ist
kein Formalismus --- eine einzige nullbare Ganzzahlspalte macht die
Merkmalsmatrix zu einem \texttt{object}-Feld, das XGBoost ablehnt und
scikit-learn stillschweigend umwandelt. Skaliert wird im Datensatz nicht; die
Standardisierung gehört zur Modellpipeline und wird je Fold nur auf den
Trainingsdaten angepasst.

Die Messlatte ist zweistufig angelegt (Abschnitt~\ref{sec:zielgroessen}).
Stufe~1 kommt ohne ein einziges Merkmal aus: der Gesamtmittelwert der
Trainingsstadtteile für die Mengenzielgrößen, die häufigste Klasse für die
Einsatzart. Sie beantwortet, ob in den Merkmalen überhaupt Information steckt.
Stufe~2 ist die einfachste Form, die zur Datenform passt: ein verallgemeinertes
lineares Modell mit kanonischem Link, unpenalisiert angepasst. Beide Stufen
laufen über dieselben 50 Läufe, dieselbe Aufteilung und dieselben zwölf
Merkmale wie die späteren Vergleichsverfahren; die Baseline ist damit
Mitbewerber unter identischem Protokoll und nicht bloß ein Erwartungswert.
Maßgeblich ist Stufe~2 --- ein Verfahren, das nur Stufe~1 schlägt, hat nichts
gezeigt (Tabelle~\ref{tab:baselines}).

Für die Mengenzielgrößen ist Stufe~2 ein Poisson-Modell mit Offset: Die
logarithmierte Wohnbevölkerung geht mit fest auf eins gesetztem Koeffizienten
ein, geschätzt wird damit die Rate (Quelltext~\ref{lst:poisson}). Über den
Logarithmus als Link wirkt das Modell auf der Skala der Zielgröße
multiplikativ; es erfasst Krümmung, aber keine Wechselwirkungen --- und genau
die finden Random Forest und XGBoost konstruktionsbedingt. Einen frei wählbaren
Parameter hat es nicht. Die zweite Mengenzielgröße entsteht aus derselben
Anpassung, geteilt durch die Wohnbevölkerung; ein eigenes Ratenmodell wäre eine
zweite Spezifikation und damit nicht mehr dasselbe Protokoll. Für die
Einsatzart ist Stufe~2 das Gegenstück: ein multinomiales Logit, ebenfalls ohne
Strafterm, mit ausgeglichenen Klassengewichten statt Resampling --- keine Zeile
wird vervielfacht oder entfernt (Quelltext~\ref{lst:logit}).

Die Zähldaten sind deutlich überdispers; der Dispersionsindex liegt bei 62,8
(Abschnitt~\ref{sec:eda}). Die naheliegende Erweiterung wäre eine negative
Binomialverteilung, und sie war ursprünglich vorgesehen. Sie löst jedoch ein
Problem, das diese Baseline nicht hat: Überdispersion beschädigt beim
Poisson-Schätzer die Standardfehler, nicht die Konsistenz des bedingten
Mittelwerts,\footcite[5]{Gourieroux1981} und eine Referenz mit reinen
Punktvorhersagen verwendet keine Standardfehler. Die Reduktion auf die
einfacheren Varianten ist damit methodisch sauber, sie vermeidet willkürliche
Parameter und liefert zudem stärkere Vergleichswerte. Tragend ist das zweite
Argument: Ein Strafterm auf dem Vorgabewert der verwendeten Software wäre ein
willkürlich gesetzter Parameter, ein optimierter eine zusätzlich zu begründende
Wahl. Das dritte gilt nicht durchgehend --- im Mengenstrang steigt die
Messlatte von 37,27 auf 33,98 RMSE, im Strukturstrang sinkt sie von 0,314 auf
0,297 Macro-F1. Dass dieselbe Änderung den einen Strang härter und den anderen
weicher macht, ist zugleich das Argument dagegen, sie nach dem Ergebnis
auszuwählen.

Drei Eigenschaften dieser Referenzwerte werden leicht als Fehler gelesen.
Negative Bestimmtheitsmaße sind hier richtig und aussagekräftig: Wer für einen
unbekannten Stadtteil den Gesamtdurchschnitt vorhersagt, liegt schlechter als
dessen eigener Mittelwert --- genau diese Lücke sollen die Strukturmerkmale
schließen. Auf der Rate ist das Bestimmtheitsmaß kein geeignetes Hauptmaß, weil
es gegen den Mittelwert der jeweiligen Testdaten misst und die Rate zwischen
den Stadtteilen um ein Vielfaches streut; maßgeblich sind dort RMSE und
mittlerer absoluter Fehler (Abschnitt~\ref{sec:modellbewertung}). Und in der
Klassifikation liegt die Trefferquote der Stufe~2 unter der von Stufe~1,
während das Macro-F1 steigt; maßgeblich ist Macro-F1, weil die Trefferquote von
der häufigsten Klasse bestimmt wird. Eine Fortschreibung des Vormonatswerts
fehlt bewusst: Sie verwendete die eigene Vergangenheit des Teststadtteils, und
ein zurückgehaltener Stadtteil hat im Training keine.

Tabelle~\ref{tab:steckbrief} fasst zusammen, was die Modellierung übernimmt.
Die letzte Zeile ist zugleich der Grund, warum die Angaben dieses Kapitels
nachprüfbar sind: 19 automatisierte Prüfungen laufen gegen die erzeugten
Dateien und nicht gegen den Code, der sie erzeugt hat.
